% Options for packages loaded elsewhere
\PassOptionsToPackage{unicode}{hyperref}
\PassOptionsToPackage{hyphens}{url}
\PassOptionsToPackage{dvipsnames,svgnames,x11names}{xcolor}
%
\documentclass[
  letterpaper,
  DIV=11,
  numbers=noendperiod]{scrartcl}

\usepackage{amsmath,amssymb}
\usepackage{iftex}
\ifPDFTeX
  \usepackage[T1]{fontenc}
  \usepackage[utf8]{inputenc}
  \usepackage{textcomp} % provide euro and other symbols
\else % if luatex or xetex
  \usepackage{unicode-math}
  \defaultfontfeatures{Scale=MatchLowercase}
  \defaultfontfeatures[\rmfamily]{Ligatures=TeX,Scale=1}
\fi
\usepackage{lmodern}
\ifPDFTeX\else  
    % xetex/luatex font selection
\fi
% Use upquote if available, for straight quotes in verbatim environments
\IfFileExists{upquote.sty}{\usepackage{upquote}}{}
\IfFileExists{microtype.sty}{% use microtype if available
  \usepackage[]{microtype}
  \UseMicrotypeSet[protrusion]{basicmath} % disable protrusion for tt fonts
}{}
\makeatletter
\@ifundefined{KOMAClassName}{% if non-KOMA class
  \IfFileExists{parskip.sty}{%
    \usepackage{parskip}
  }{% else
    \setlength{\parindent}{0pt}
    \setlength{\parskip}{6pt plus 2pt minus 1pt}}
}{% if KOMA class
  \KOMAoptions{parskip=half}}
\makeatother
\usepackage{xcolor}
\setlength{\emergencystretch}{3em} % prevent overfull lines
\setcounter{secnumdepth}{-\maxdimen} % remove section numbering
% Make \paragraph and \subparagraph free-standing
\makeatletter
\ifx\paragraph\undefined\else
  \let\oldparagraph\paragraph
  \renewcommand{\paragraph}{
    \@ifstar
      \xxxParagraphStar
      \xxxParagraphNoStar
  }
  \newcommand{\xxxParagraphStar}[1]{\oldparagraph*{#1}\mbox{}}
  \newcommand{\xxxParagraphNoStar}[1]{\oldparagraph{#1}\mbox{}}
\fi
\ifx\subparagraph\undefined\else
  \let\oldsubparagraph\subparagraph
  \renewcommand{\subparagraph}{
    \@ifstar
      \xxxSubParagraphStar
      \xxxSubParagraphNoStar
  }
  \newcommand{\xxxSubParagraphStar}[1]{\oldsubparagraph*{#1}\mbox{}}
  \newcommand{\xxxSubParagraphNoStar}[1]{\oldsubparagraph{#1}\mbox{}}
\fi
\makeatother

\usepackage{color}
\usepackage{fancyvrb}
\newcommand{\VerbBar}{|}
\newcommand{\VERB}{\Verb[commandchars=\\\{\}]}
\DefineVerbatimEnvironment{Highlighting}{Verbatim}{commandchars=\\\{\}}
% Add ',fontsize=\small' for more characters per line
\usepackage{framed}
\definecolor{shadecolor}{RGB}{241,243,245}
\newenvironment{Shaded}{\begin{snugshade}}{\end{snugshade}}
\newcommand{\AlertTok}[1]{\textcolor[rgb]{0.68,0.00,0.00}{#1}}
\newcommand{\AnnotationTok}[1]{\textcolor[rgb]{0.37,0.37,0.37}{#1}}
\newcommand{\AttributeTok}[1]{\textcolor[rgb]{0.40,0.45,0.13}{#1}}
\newcommand{\BaseNTok}[1]{\textcolor[rgb]{0.68,0.00,0.00}{#1}}
\newcommand{\BuiltInTok}[1]{\textcolor[rgb]{0.00,0.23,0.31}{#1}}
\newcommand{\CharTok}[1]{\textcolor[rgb]{0.13,0.47,0.30}{#1}}
\newcommand{\CommentTok}[1]{\textcolor[rgb]{0.37,0.37,0.37}{#1}}
\newcommand{\CommentVarTok}[1]{\textcolor[rgb]{0.37,0.37,0.37}{\textit{#1}}}
\newcommand{\ConstantTok}[1]{\textcolor[rgb]{0.56,0.35,0.01}{#1}}
\newcommand{\ControlFlowTok}[1]{\textcolor[rgb]{0.00,0.23,0.31}{\textbf{#1}}}
\newcommand{\DataTypeTok}[1]{\textcolor[rgb]{0.68,0.00,0.00}{#1}}
\newcommand{\DecValTok}[1]{\textcolor[rgb]{0.68,0.00,0.00}{#1}}
\newcommand{\DocumentationTok}[1]{\textcolor[rgb]{0.37,0.37,0.37}{\textit{#1}}}
\newcommand{\ErrorTok}[1]{\textcolor[rgb]{0.68,0.00,0.00}{#1}}
\newcommand{\ExtensionTok}[1]{\textcolor[rgb]{0.00,0.23,0.31}{#1}}
\newcommand{\FloatTok}[1]{\textcolor[rgb]{0.68,0.00,0.00}{#1}}
\newcommand{\FunctionTok}[1]{\textcolor[rgb]{0.28,0.35,0.67}{#1}}
\newcommand{\ImportTok}[1]{\textcolor[rgb]{0.00,0.46,0.62}{#1}}
\newcommand{\InformationTok}[1]{\textcolor[rgb]{0.37,0.37,0.37}{#1}}
\newcommand{\KeywordTok}[1]{\textcolor[rgb]{0.00,0.23,0.31}{\textbf{#1}}}
\newcommand{\NormalTok}[1]{\textcolor[rgb]{0.00,0.23,0.31}{#1}}
\newcommand{\OperatorTok}[1]{\textcolor[rgb]{0.37,0.37,0.37}{#1}}
\newcommand{\OtherTok}[1]{\textcolor[rgb]{0.00,0.23,0.31}{#1}}
\newcommand{\PreprocessorTok}[1]{\textcolor[rgb]{0.68,0.00,0.00}{#1}}
\newcommand{\RegionMarkerTok}[1]{\textcolor[rgb]{0.00,0.23,0.31}{#1}}
\newcommand{\SpecialCharTok}[1]{\textcolor[rgb]{0.37,0.37,0.37}{#1}}
\newcommand{\SpecialStringTok}[1]{\textcolor[rgb]{0.13,0.47,0.30}{#1}}
\newcommand{\StringTok}[1]{\textcolor[rgb]{0.13,0.47,0.30}{#1}}
\newcommand{\VariableTok}[1]{\textcolor[rgb]{0.07,0.07,0.07}{#1}}
\newcommand{\VerbatimStringTok}[1]{\textcolor[rgb]{0.13,0.47,0.30}{#1}}
\newcommand{\WarningTok}[1]{\textcolor[rgb]{0.37,0.37,0.37}{\textit{#1}}}

\providecommand{\tightlist}{%
  \setlength{\itemsep}{0pt}\setlength{\parskip}{0pt}}\usepackage{longtable,booktabs,array}
\usepackage{calc} % for calculating minipage widths
% Correct order of tables after \paragraph or \subparagraph
\usepackage{etoolbox}
\makeatletter
\patchcmd\longtable{\par}{\if@noskipsec\mbox{}\fi\par}{}{}
\makeatother
% Allow footnotes in longtable head/foot
\IfFileExists{footnotehyper.sty}{\usepackage{footnotehyper}}{\usepackage{footnote}}
\makesavenoteenv{longtable}
\usepackage{graphicx}
\makeatletter
\newsavebox\pandoc@box
\newcommand*\pandocbounded[1]{% scales image to fit in text height/width
  \sbox\pandoc@box{#1}%
  \Gscale@div\@tempa{\textheight}{\dimexpr\ht\pandoc@box+\dp\pandoc@box\relax}%
  \Gscale@div\@tempb{\linewidth}{\wd\pandoc@box}%
  \ifdim\@tempb\p@<\@tempa\p@\let\@tempa\@tempb\fi% select the smaller of both
  \ifdim\@tempa\p@<\p@\scalebox{\@tempa}{\usebox\pandoc@box}%
  \else\usebox{\pandoc@box}%
  \fi%
}
% Set default figure placement to htbp
\def\fps@figure{htbp}
\makeatother

\usepackage{float}
\usepackage{tabularray}
\usepackage[normalem]{ulem}
\usepackage{graphicx}
\UseTblrLibrary{booktabs}
\UseTblrLibrary{rotating}
\UseTblrLibrary{siunitx}
\NewTableCommand{\tinytableDefineColor}[3]{\definecolor{#1}{#2}{#3}}
\newcommand{\tinytableTabularrayUnderline}[1]{\underline{#1}}
\newcommand{\tinytableTabularrayStrikeout}[1]{\sout{#1}}
\usepackage{fvextra}
\DefineVerbatimEnvironment{Highlighting}{Verbatim}{breaklines,commandchars=\\\{\}}
\DefineVerbatimEnvironment{OutputCode}{Verbatim}{breaklines,commandchars=\\\{\}}
\KOMAoption{captions}{tableheading}
\makeatletter
\@ifpackageloaded{caption}{}{\usepackage{caption}}
\AtBeginDocument{%
\ifdefined\contentsname
  \renewcommand*\contentsname{Table of contents}
\else
  \newcommand\contentsname{Table of contents}
\fi
\ifdefined\listfigurename
  \renewcommand*\listfigurename{List of Figures}
\else
  \newcommand\listfigurename{List of Figures}
\fi
\ifdefined\listtablename
  \renewcommand*\listtablename{List of Tables}
\else
  \newcommand\listtablename{List of Tables}
\fi
\ifdefined\figurename
  \renewcommand*\figurename{Figure}
\else
  \newcommand\figurename{Figure}
\fi
\ifdefined\tablename
  \renewcommand*\tablename{Table}
\else
  \newcommand\tablename{Table}
\fi
}
\@ifpackageloaded{float}{}{\usepackage{float}}
\floatstyle{ruled}
\@ifundefined{c@chapter}{\newfloat{codelisting}{h}{lop}}{\newfloat{codelisting}{h}{lop}[chapter]}
\floatname{codelisting}{Listing}
\newcommand*\listoflistings{\listof{codelisting}{List of Listings}}
\makeatother
\makeatletter
\makeatother
\makeatletter
\@ifpackageloaded{caption}{}{\usepackage{caption}}
\@ifpackageloaded{subcaption}{}{\usepackage{subcaption}}
\makeatother

\usepackage{bookmark}

\IfFileExists{xurl.sty}{\usepackage{xurl}}{} % add URL line breaks if available
\urlstyle{same} % disable monospaced font for URLs
\hypersetup{
  pdftitle={report},
  pdfauthor={Hyoungchul Kim},
  colorlinks=true,
  linkcolor={blue},
  filecolor={Maroon},
  citecolor={Blue},
  urlcolor={Blue},
  pdfcreator={LaTeX via pandoc}}


\title{report}
\author{Hyoungchul Kim}
\date{2025-03-21}

\begin{document}
\maketitle


\section{Lorem ipsum}\label{lorem-ipsum}

\subsection{Bullet point}\label{bullet-point}

This is a template for my report.

\begin{enumerate}
\def\labelenumi{\arabic{enumi}.}
\tightlist
\item
  Lorem.
\end{enumerate}

\begin{itemize}
\tightlist
\item
  Lorem.
\end{itemize}

\subsection{Math equations}\label{math-equations}

This is a latex math-equation using \emph{align}.

\begin{align*}
  (1) + (2) + (3) \geq 3
\end{align*}

--

\subsection{Code block}\label{code-block}

This is a code block.

\begin{Shaded}
\begin{Highlighting}[]
\FunctionTok{library}\NormalTok{(fixest)}
\NormalTok{  mods }\OtherTok{=} \FunctionTok{feols}\NormalTok{(}
\NormalTok{  rating }\SpecialCharTok{\textasciitilde{}}\NormalTok{ complaints }\SpecialCharTok{+}\NormalTok{ privileges }\SpecialCharTok{+}\NormalTok{ learning }\SpecialCharTok{+} \FunctionTok{csw0}\NormalTok{(raises }\SpecialCharTok{+}\NormalTok{ critical), }\AttributeTok{data =}\NormalTok{ attitude}
\NormalTok{)}

\NormalTok{dict }\OtherTok{=} \FunctionTok{c}\NormalTok{(}\StringTok{"rating"}     \OtherTok{=} \StringTok{"Overall Rating"}\NormalTok{,}
         \StringTok{"complaints"} \OtherTok{=} \StringTok{"Handling of Complaints"}\NormalTok{,}
         \StringTok{"privileges"} \OtherTok{=} \StringTok{"No Special Priviledges"}\NormalTok{,}
         \StringTok{"learning"}   \OtherTok{=} \StringTok{"Opportunity to Learn"}\NormalTok{,}
         \StringTok{"raises"}     \OtherTok{=} \StringTok{"Performance{-}Based Raises"}\NormalTok{,}
         \StringTok{"critical"}   \OtherTok{=} \StringTok{"Too Critical"}\NormalTok{)}

\FunctionTok{library}\NormalTok{(modelsummary) }\CommentTok{\# Make sure you have \textgreater{}=v2.0.0}

\FunctionTok{modelsummary}\NormalTok{(}
  \FunctionTok{setNames}\NormalTok{(mods, }\FunctionTok{c}\NormalTok{(}\StringTok{"(1)"}\NormalTok{, }\StringTok{"(2)"}\NormalTok{)),}
  \AttributeTok{coef\_map =}\NormalTok{ dict, }\AttributeTok{stars =} \ConstantTok{TRUE}\NormalTok{,}
  \AttributeTok{gof\_map =} \ConstantTok{NA}
\NormalTok{  ) }\SpecialCharTok{|\textgreater{}}
  \CommentTok{\# some optional stylistic tweaks}
\NormalTok{  tinytable}\SpecialCharTok{::}\FunctionTok{group\_tt}\NormalTok{(}\AttributeTok{j =} \FunctionTok{list}\NormalTok{(}\StringTok{"Dep. variable: Overall Rating"} \OtherTok{=} \DecValTok{2}\SpecialCharTok{:}\DecValTok{3}\NormalTok{)) }\SpecialCharTok{|\textgreater{}}
\NormalTok{  tinytable}\SpecialCharTok{::}\FunctionTok{style\_tt}\NormalTok{(}\AttributeTok{i =} \DecValTok{1}\SpecialCharTok{:}\DecValTok{2}\NormalTok{, }\AttributeTok{j =} \DecValTok{2}\SpecialCharTok{:}\DecValTok{3}\NormalTok{, }\AttributeTok{background =} \StringTok{"pink"}\NormalTok{)}
\end{Highlighting}
\end{Shaded}

\begin{table}
\centering
\begin{talltblr}[         %% tabularray outer open
entry=none,label=none,
note{}={+ p \num{< 0.1}, * p \num{< 0.05}, ** p \num{< 0.01}, *** p \num{< 0.001}},
]                     %% tabularray outer close
{                     %% tabularray inner open
colspec={Q[]Q[]Q[]},
column{1}={}{halign=l,},
cell{1}{3}={}{halign=c,},
cell{2}{2}={}{halign=c,},
cell{2}{3}={}{halign=c,},
cell{5}{2}={}{halign=c,},
cell{5}{3}={}{halign=c,},
cell{6}{2}={}{halign=c,},
cell{6}{3}={}{halign=c,},
cell{7}{2}={}{halign=c,},
cell{7}{3}={}{halign=c,},
cell{8}{2}={}{halign=c,},
cell{8}{3}={}{halign=c,},
cell{9}{2}={}{halign=c,},
cell{9}{3}={}{halign=c,},
cell{10}{2}={}{halign=c,},
cell{10}{3}={}{halign=c,},
cell{11}{2}={}{halign=c,},
cell{11}{3}={}{halign=c,},
cell{12}{2}={}{halign=c,},
cell{12}{3}={}{halign=c,},
cell{3}{2}={}{halign=c, bg=pink,},
cell{3}{3}={}{halign=c, bg=pink,},
cell{4}{2}={}{halign=c, bg=pink,},
cell{4}{3}={}{halign=c, bg=pink,},
cell{1}{2}={c=2,}{halign=c, halign=c,},
}                     %% tabularray inner close
\toprule
& Dep. variable: Overall Rating &  \\ \cmidrule[lr]{2-3}
& (1) & (2) \\ \midrule %% TinyTableHeader
Handling of Complaints   & \num{0.682}*** & \num{0.692}*** \\
& (\num{0.129})  & (\num{0.149})  \\
No Special Priviledges   & \num{-0.103}   & \num{-0.104}   \\
& (\num{0.129})  & (\num{0.135})  \\
Opportunity to Learn     & \num{0.238}+   & \num{0.249}    \\
& (\num{0.139})  & (\num{0.160})  \\
Performance-Based Raises &                 & \num{-0.033}   \\
&                 & (\num{0.202})  \\
Too Critical             &                 & \num{0.015}    \\
&                 & (\num{0.147})  \\
\bottomrule
\end{talltblr}
\end{table}

\begin{Shaded}
\begin{Highlighting}[]
\CommentTok{\# You need reticulate R{-}package for this}
\BuiltInTok{print}\NormalTok{(}\StringTok{"hello, python world."}\NormalTok{)}
\end{Highlighting}
\end{Shaded}

\begin{verbatim}
hello, python world.
\end{verbatim}

\begin{Shaded}
\begin{Highlighting}[]

\CommentTok{\# You need JuliaCall R{-}package for this}
\FunctionTok{println}\NormalTok{(}\StringTok{"hello, julia world."}\NormalTok{)}
\end{Highlighting}
\end{Shaded}

\begin{verbatim}
hello, julia world.
\end{verbatim}




\end{document}
